\section{Geospatial grounding limits for scientific evidence}

\subsection{Benchmark methodology}

Gemini 3.5 Flash (Api-Revision 2026-05-20) with the \texttt{google\_maps} tool at San Juan Islands coordinates (48.5465, $-$123.03) was measured across three query types: total citation count, citation type split (place vs.\ scientific/dataset), grounding density per thousand words, and the unsupported scientific claim rate:

\equncited The orcast comparison used the same orca evidence query submitted to \texttt{/api/interactions/plan}, counting annotations with a bound artifact or href. The measurement tool is available at \texttt{tools/testing/maps\_grounding\_probe.py}.

\subsection{Results}

Across three Maps grounding queries, 0 of 25 citations resolved to scientific or dataset sources; all 25 resolved to Google Maps place pins. The evidence query produced 9 scientific sentences (NOAA fisheries data, Center for Whale Research census, Chinook salmon corridors, J/K/L pod hydrophone data); 8 of 9 (89\%) were uncited. Averaged across evidence-bearing queries: $R_\mathrm{uncited} = 85\%$. The orcast result: 6 of 6 annotations were bound; $R_\mathrm{uncited} = 0\%$.

\subsection{Interpretation}

This reflects geospatial grounding's structural scope boundary, not a Gemini or Maps deficiency. Maps grounding is designed for place information and excels at that: the POI query achieved 8 citations in 229 words with 0 unsupported sentences. Place grounding cannot resolve scientific evidence claims to their origin data, method, or experiment. The LLM geospatial knowledge study~\cite{arxiv231013002} provides peer-reviewed evidence that LLM geospatial competence degrades substantially for knowledge not encoded as named places. A domain-specific provenance system is required for the evidence layer.

\subsection{Reproducibility}

Benchmark run: 2026-06-24. Model: gemini-3.5-flash. API revision: 2026-05-20. Coordinates: 48.5465, $-$123.03. Fully reproducible via \texttt{tools/testing/maps\_grounding\_probe.py} with a Gemini API key.

\subsection{Extended benchmark: 8-scenario parallel RAG comparison}

A follow-on 8-scenario parallel benchmark (\texttt{tools/testing/grounding\_parallel\_rag.py}, 2026-06-24) compared Maps-only grounding against five RAG-augmented variants using live orcast skill outputs as context. The hierarchy of $R_\mathrm{uncited}$ across conditions reveals a diagnostic structure:

\begin{center}
\small
\begin{tabular}{@{}lll@{}}
  \toprule
  Scenario & Context injected & $R_\mathrm{uncited}$ \\
  \midrule
  S4 (surface planner step-log) & Orchestrated reasoning trace & \textbf{0\%} \\
  S8 (hotspot data) & Structured probability data & 75\% \\
  S1, S7 (Maps-only baseline) & None & 60--100\% \\
  S2 (gate context) & Narrow gate data & 100\% \\
  \bottomrule
\end{tabular}
\end{center}

The surface planner step-log (Scenario~4) eliminates $R_\mathrm{uncited}$ entirely by transforming the query from an open-domain science query into a closed artifact-reference query: the model answers from the step-log rather than from world knowledge. This is not improved citation generation; it is query transformation. Structured probability data alone (Scenario~8, hotspot injection) lowers $R_\mathrm{uncited}$ to 75\%, a 25 percentage point reduction from the Scenario~7 Maps-only baseline (100\%), but does not eliminate uncited claims. Gate context (Scenario~2) does not reduce $R_\mathrm{uncited}$ because the injected gate data does not cover the scientific claims Gemini independently generates. The hierarchy functions as a diagnostic for grounding architecture quality: context that binds the query to a complete reasoning trace eliminates uncited claims; narrowly structured context does not. This finding motivates a formal evaluation framework for AI reasoning systems (see companion paper, \textit{Grounding quality measurement for orchestrated AI reasoning chains})~\cite{arxiv231109533,arxiv240115884}.
